\newpage
\section{Vektoren und Punkte}

\subsection{Unterschiede}
Vektoren und Punkte sind unterschiedliche Objekte mit unterschiedlichen Eigenschaften.

Ein \textbf{Punkt} befindet sich an einem bestimmten Ort und er hat keine Richtung. Ein Punkt $P$ mit den \textbf{Koordinaten} $x$ und $y$ wird wie folgt geschrieben:
\[
  \coord[P]{x,y}
\]
Ein \textbf{Vektor} kann beliebig verschoben werden, er hat also keine feste Position, aber er zeigt in eine bestimmte Richtung. Ein Vektor $\vv{v}$ mit den \textbf{Komponenten} $v_{x}$ und $v_{y}$ wird folgendermassen geschrieben:
\[
  \vv{v} = \dvec{v_{x},v_{y}}
\]

% ------------------------------------------------------------------------------
\subsection{Ortsvektor}

Wird ein Vektor $\vv{OP}$ mit dem Ursprung $O$ als Anfangspunkt und dem Punkt $\coord[P]{x,y}$ als Endpunkt definiert wird, so entsprechen die Komponenten des Vektors gerade den Koordinaten des Punkts $P$:
\[
  \vv{OP} = \dvec{x-0,y-0} = \dvec{x,y}
\]
\begin{center}
  \begin{tikzpicture}
    \tkzInit[xmin=-0.5,xmax=7,ymin=-0.5,ymax=3]
    \tkzDrawXY
    \tkzDefPoint(0,0){O}
    \tkzDefPoint(6,2){P}

    \tkzDrawSegment[dashed,red,thick,vector style](O,P)
    \tkzLabelSegment[red,above left](O,P){$\vv{OP}=\tvec{x,y}$}

    \tkzDrawPoints(O,P)
    \tkzLabelPoint[above left](O){$O$}
    \tkzLabelPoint[above right](P){$P\coord{x,y}$}
  \end{tikzpicture}
\end{center}

Dieser Vektor wird \textbf{Ortsvektor} des Punkts $P$ genannt. Der Ortsvektor von $P$ zeigt vom Ursprung $O$ zum Punkt $P$.

Werden Vektoren als Verschiebungen aufgefasst, so verschiebt der Ortsvektor den Ursprung $O$ in den Punkt $P$.
\begin{example}
  \textbf{Beispiel:} Der Ortsvektor des Punkts $\coord[A]{-3,5}$ lautet $\vv{OA} = \tvec{-3,5}$.
\end{example}

% ------------------------------------------------------------------------------
\subsection{Vektor zwischen Punkten}

Ein Vektor kann durch einen Anfangspunkt $A$ und einen Endpunkt $B$ bestimmt werden.
\begin{center}
  \begin{tikzpicture}
    \tkzInit[xmin=-0.5,xmax=10,ymin=-0.5,ymax=4]
    \tkzDrawXY
    \tkzDefPoint(0,0){O}
    \tkzDefPoint(2,3){A}
    \tkzDefPoint(9,1){B}

    \tkzDrawSegment[dashed,vector style](O,A)
    \tkzLabelSegment[above left](O,A){$\vv{OA}$}

    \tkzDrawSegment[dashed,vector style](O,B)
    \tkzLabelSegment[above left](O,B){$\vv{OB}$}

    \tkzDrawSegment[red,thick,vector style](A,B)
    \tkzLabelSegment[red,above right](A,B){$\vv{AB}$}

    \tkzDrawPoints(A,B)
    \tkzLabelPoints[above left](O)
    \tkzLabelPoint[above](A){$A\coord{x_{A},y_{A}}$}
    \tkzLabelPoint[right](B){$B\coord{x_{B},y_{B}}$}
  \end{tikzpicture}
\end{center}
Der Vektor $\vv{AB}$, welcher von Punkt $A$ nach Punkt $B$ zeigt, kann als Subtraktion der Ortsvektoren von $A$ und $B$ ausgedrückt werden.
\[
  \vv{AB} = \vv{OB} - \vv{OA} = \dvec{x_{B},y_{B}} - \dvec{x_{A},y_{A}} = \dvec{x_{B}-x_{A},y_{B}-y_{A}}
\]
Daraus lässt sich schliessen, dass sich die Komponenten des Vektors $\vv{AB}$ durch Subtraktion der Koordinaten von Punkt $B$ und Punkt $A$ ergeben.
\begin{theorem}
  \textbf{Vektor zwischen Punkten.} Hat ein Vektor $\vv{AB}$ den Anfangspunkt $\coord[A]{x_{A},y_{A}}$ und den Endpunkt $\coord[B]{x_{B},y_{B}}$ besitzt, so lautet seine Komponentendarstellung
  \[
    \vv{AB} = \dvec{x_{B}-x_{A},y_{B}-y_{A}}
  \]
\end{theorem}
\begin{example}
  \textbf{Beispiel:} Im der oben stehenden Abbildung ist der Vektor $\vv{v}$ durch den Anfangspunkt $\coord[A]{2,3}$ und den Endpunkt $\coord[B]{9,1}$ definiert. Also lautet seine Komponentendarstellung
  \[
    \vv{AB} = \dvec{9-2,1-3} = \dvec{7,-2}
  \]
\end{example}
